\documentclass[11pt]{amsart}

\usepackage[T1]{fontenc}
\usepackage{newtxtext,newtxmath}
\usepackage[margin=1.25in]{geometry}
\usepackage{hyperref}

\title[Discovering Hidden Relations Among Triangle Invariants]%
{Discovering Hidden Relations \\ Among Triangle Invariants}

\author{triangle-relations project}
\date{}

\subjclass[2020]{51M04, 68T07}
\keywords{triangle geometry, Euler's relation, functional dependence, autoencoders, dimension counting}

\newcommand{\M}{\mathcal{M}}

\newtheorem{proposition}{Proposition}
\newtheorem{remark}[proposition]{Remark}

\begin{document}

\begin{abstract}
We describe a numerical method for discovering previously unknown functional
relations among scalar quantities derived from a triangle, in the spirit of
Euler's classical relation between the circumradius, the inradius, and the
distance between the incenter and circumcenter. Because a triangle has three
degrees of freedom up to rigid motion, a relation among four or more derived
scalars is guaranteed by dimension counting alone and is therefore
uninformative; a relation among exactly three scalars, by contrast, is a
genuine and generically unexpected algebraic coincidence. We detect such coincidences by training a bottleneck autoencoder on sampled
triangle data and comparing its reconstruction error against a
column-permuted null. Applied to the circumradius, inradius, and
incenter-to-circumcenter distance, the method flags exactly this triple as
functionally dependent, with no prior knowledge of Euler's relation.
\end{abstract}

\maketitle

\section{Setup: the moduli space of triangles}

A triangle is specified by three points in $\mathbb{R}^2$, i.e.\ six real
parameters. Every quantity we consider (area, perimeter, circumradius $R$,
inradius $r$, distances between derived points, \dots) is built purely from
the vertex coordinates and is invariant under rigid motions: translating or
rotating the triangle does not change its area, its inradius, or the
distance between its incenter and circumcenter. Translations (2 parameters)
and rotations (1 parameter) form a 3-dimensional group acting on the
6-dimensional space of vertex triples; quotienting it out leaves a
$3$-dimensional moduli space $\M$ of triangle shapes up to congruence
(reflection is a further discrete symmetry and does not change this
dimension count). Concretely, $\M$ can be coordinatized by, e.g., the three
side lengths $(a,b,c)$ subject to the triangle inequality.

Every derived scalar quantity is a smooth function $f : \M \to \mathbb{R}$.
Choosing $k$ such quantities gives a smooth map
\[
  F = (f_1, \dots, f_k) : \M \to \mathbb{R}^k .
\]
Recall that a (smooth) $m$-dimensional submanifold of $\mathbb R^n$ is a
subset that, near each of its points, agrees with the graph of a smooth
function expressing $n-m$ of the coordinates in terms of the other $m$; this
is what lets us speak of $\M$ itself, and of the image of $F$, as having a
well-defined dimension even though neither is simply a coordinate subspace
of $\mathbb R^k$. We call a property of such a map \emph{generic} if it
holds for every choice of $F$ outside a closed, measure-zero exceptional
set — informally, for ``almost every'' choice of derived quantities, and
stably so under small perturbations of them.

\section{Why relations among four or more quantities are not interesting}

If $k > \dim \M = 3$, a generic smooth map from a 3-manifold into
$\mathbb{R}^k$ has, as its image, an (at most) 3-dimensional submanifold of
$\mathbb{R}^k$ — this is just the statement that embedding a 3-dimensional
space into a higher-dimensional ambient space produces $k-3$ constraint
equations for free, regardless of which $k$ functions were chosen, as long
as they are not specially degenerate. Concretely, this image is cut out
locally by $k-3$ independent equations $G_i(f_1,\dots,f_k) = 0$. So finding
\emph{some} relation among four or more triangle invariants is guaranteed by
dimension counting alone. It says nothing specific about triangle geometry —
it would happen for essentially any four smooth functions on any
3-dimensional space of objects.

\section{Why a relation among exactly three quantities is special}

If $k = \dim \M = 3$ exactly, the situation is qualitatively different.
Recall that the Jacobian $DF(p)$ of a differentiable map
$F=(f_1,f_2,f_3):\M\to\mathbb R^3$ at a point $p$ is the $3\times 3$ matrix
of partial derivatives $\big(\partial f_i/\partial x_j\big)$ in local
coordinates $x_j$ on $\M$; its rank measures how many of the $f_i$ vary
independently of one another near $p$. Recall also that $F$ is a \emph{local
diffeomorphism} at $p$ if it restricts to a smooth bijection, with smooth
inverse, from some neighborhood of $p$ onto some neighborhood of $F(p)$ —
by the inverse function theorem, exactly when $DF(p)$ is invertible.

A generic smooth map $F : \M \to \mathbb{R}^3$ is, at a generic point, a
local diffeomorphism: its Jacobian $DF$ has full rank 3, and as the triangle
shape ranges over an open subset of $\M$, the triple $(f_1,f_2,f_3)$ sweeps
out a full 3-dimensional \emph{open region} of $\mathbb{R}^3$ — a solid
blob, not a surface. This is again a genericity statement: for an arbitrary
hand-picked triple of triangle invariants, there is no a priori reason for
$DF$ to be rank-deficient anywhere.

An \emph{exact} relation $G(f_1,f_2,f_3) = 0$ holding on all of $\M$ is
therefore equivalent to $DF$ having rank $\le 2$ everywhere: the image
collapses onto a 2-dimensional surface embedded in $\mathbb{R}^3$, instead of
filling out a full 3-dimensional region. That is not generic — it is a
genuine algebraic coincidence, exactly the kind of fact that constitutes a
classical theorem. Euler's relation
\[
  d^2 = R^2 - 2Rr
\]
(where $d$ is the distance between the incenter and circumcenter) is exactly
a statement that the Jacobian of $(R, r, d)$ as functions of triangle shape
has rank $\le 2$ everywhere, collapsing three ostensibly independent
quantities onto a single surface. This is the precise sense in which
searching for relations among exactly three derived scalars — and not four —
is the search for something structurally new.

\section{Detecting rank-deficiency with a bottleneck autoencoder}

Given a candidate triple $(f_1,f_2,f_3)$, we want a numerical test for
whether the induced map $F : \M \to \mathbb{R}^3$ has full rank 3 generically
(data fills a 3D region) or is everywhere rank $\le 2$ (data is confined to a
2D surface). We sample many random triangles — many points of $\M$ — and
evaluate $(f_1,f_2,f_3)$ on each, producing a point cloud
$X \subset \mathbb{R}^3$.

An autoencoder with a 2-dimensional latent bottleneck, an encoder
$e : \mathbb{R}^3 \to \mathbb{R}^2$ and decoder $d : \mathbb{R}^2 \to
\mathbb{R}^3$ trained to minimize $\mathbb{E}\, \| d(e(x)) - x \|^2$ over the
sample, is exactly searching for the best rank-2 (nonlinear) approximation to
the data:
\begin{itemize}
  \item If the data truly lies on a 2-dimensional surface (an exact relation
    holds), a 2-dimensional latent code is not actually a bottleneck for it —
    a sufficiently expressive encoder/decoder pair can in principle drive the
    reconstruction error to (near) zero, since the encoder just needs to
    learn a local chart of the surface and the decoder its inverse.
  \item If the data genuinely fills an open 3-dimensional region, no
    continuous map that factors through a 2-dimensional intermediate space
    can reconstruct it without loss: an entire dimension worth of variation
    is necessarily discarded. Reconstruction error is then bounded away from
    zero, on the order of the data's own extent in the discarded direction.
\end{itemize}

This gives a clean signature: near-zero reconstruction error with a
2-dimensional bottleneck is evidence for an (exact or near-exact) relation
among the three chosen quantities; error comparable to the data's own scale
is consistent with the generic, full-rank, no-relation case.

\subsection{A precise statement}

We can make this dichotomy precise. Let $P$ denote the law of
$(f_1,f_2,f_3)(T)$ for a random triangle $T$, and write
$X = \operatorname{supp}(P) \subset \mathbb{R}^3$ for the support of $P$
(the smallest closed set carrying all of its probability mass).

By a \emph{two-layer perceptron of hidden width $h$} we mean a function
$N : \mathbb R^p \to \mathbb R^q$ of the form
\[
  N(x) \;=\; W_2 \, \sigma(W_1 x + b_1) + b_2 ,
\]
where $W_1 \in \mathbb R^{h \times p}$, $b_1 \in \mathbb R^h$,
$W_2 \in \mathbb R^{q \times h}$, $b_2 \in \mathbb R^q$ are the trainable
weights and biases (``two-layer'' for the two weight matrices $W_1, W_2$;
this is what our encoders and decoders are individually, before being
composed into the single autoencoder network of \S 4), and $\sigma$ is a
fixed nonlinear activation applied componentwise — here $\sigma = \tanh$.
Recall that a function $g$ is \emph{$K$-Lipschitz} if
$\|g(x) - g(y)\| \le K\|x - y\|$ for all $x,y$; since $\tanh$ itself is
$1$-Lipschitz, every two-layer perceptron is Lipschitz, with constant
controlled by the operator norms of $W_1$ and $W_2$. For a fixed hidden
width $h$, let $\mathcal N_h(\mathbb R^p, \mathbb R^q)$ denote the class of
all two-layer perceptrons $\mathbb R^p \to \mathbb R^q$ of that width. We
train an encoder $e \in \mathcal N_h(\mathbb R^3, \mathbb R^2)$ and a
decoder $d \in \mathcal N_h(\mathbb R^2, \mathbb R^3)$ to minimize the
reconstruction risk
\[
  L(e,d) \;=\; \mathbb{E}_{x \sim P} \, \|d(e(x)) - x\|^2 .
\]

\begin{proposition}[Sufficiency]
Recall that a subset of $\mathbb R^n$ is \emph{compact} exactly when it is
closed and bounded (Heine--Borel) — for us, $X$ is compact whenever the
random triangles are sampled from a bounded region, as they are in
practice. If $X$ is compact and admits continuous maps $\varphi : X \to \mathbb R^2$
and $\psi : \varphi(X) \to \mathbb R^3$ with $\psi(\varphi(x)) = x$ for all
$x \in X$ — in particular, whenever an exact relation
$G(f_1,f_2,f_3) = 0$ confines $X$ to a $2$-dimensional set — then for every
$\varepsilon > 0$ there is a width $h$ and a pair
$(e,d) \in \mathcal N_h(\mathbb R^3,\mathbb R^2) \times
\mathcal N_h(\mathbb R^2,\mathbb R^3)$ with $L(e,d) < \varepsilon$.
\end{proposition}

\begin{proof}[Proof sketch]
Recall the universal approximation theorem (Cybenko, 1989; Hornik, 1991):
for any continuous $\varphi : K \to \mathbb R^q$ on a compact set $K$ and
any $\delta > 0$, there is a hidden width $h$ and a two-layer perceptron
$N \in \mathcal N_h(\mathbb R^p, \mathbb R^q)$ with
$\sup_{x \in K} \|N(x) - \varphi(x)\| < \delta$. Applying this to $\varphi$
on $X$ and to $\psi$ on $\varphi(X)$ gives encoders/decoders $e, d$ with
$\sup_{x \in X} \|d(e(x)) - x\| \to 0$ as $\delta \to 0$, hence
$L(e,d) \to 0$.
\end{proof}

\begin{proposition}[Necessity, at bounded capacity]
Fix $K > 0$. Recall that a probability measure $P$ on $\mathbb R^n$ is
\emph{absolutely continuous} (with respect to Lebesgue measure) if it
assigns zero probability to every set of Lebesgue measure zero, i.e.\ it has
a density; this is the standard way of saying ``$P$ does not secretly
concentrate on some lower-dimensional set.'' Suppose $X$ contains an open
ball $B \subset \mathbb R^3$ on which $P$ is absolutely continuous with
$P(B) > 0$. Then there is a constant $c = c(P|_B, K) > 0$ such that
$L(e,d) \ge c$ for every pair of $K$-Lipschitz maps
$e : \mathbb R^3 \to \mathbb R^2$, $d : \mathbb R^2 \to \mathbb R^3$ — in
particular for every $(e,d) \in \mathcal N_h(\mathbb R^3,\mathbb R^2) \times
\mathcal N_h(\mathbb R^2,\mathbb R^3)$, of any width $h$, whose realized
weights keep the Lipschitz constant below $K$.
\end{proposition}

\begin{proof}[Proof sketch]
Recall that the Hausdorff dimension of a set $S \subset \mathbb R^n$
generalizes the usual notion of dimension to arbitrary (possibly non-smooth)
sets: informally, it is the value $\alpha$ such that covering $S$ by balls
of radius $\varepsilon$ requires on the order of $\varepsilon^{-\alpha}$ of
them as $\varepsilon \to 0$; it agrees with the usual dimension on smooth
submanifolds, and a standard fact is that Lipschitz maps do not increase
it. Since $d$ is $K$-Lipschitz, its image $d(\mathbb R^2)$ therefore has
Hausdorff dimension at most $2$, and a subset of $\mathbb R^3$ of Hausdorff
dimension less than $3$ automatically has (ordinary, $3$-dimensional)
Lebesgue measure zero, i.e.\ is \emph{Lebesgue-null}. As $P$ is absolutely
continuous on $B$ with $P(B) > 0$, it assigns zero mass to this null set, so
$\operatorname{dist}(x, d(\mathbb R^2)) > 0$ for $P|_B$-almost every $x$; a
routine compactness argument over $K$-Lipschitz maps then gives a strictly
positive lower bound $c$ on $\mathbb E_{x \sim P|_B}
\operatorname{dist}(x, d(\mathbb R^2))^2$. Since $d(e(\mathbb R^3)) \subseteq
d(\mathbb R^2)$, we have $\|d(e(x)) - x\| \ge \operatorname{dist}(x,
d(\mathbb R^2))$ pointwise, and the same bound applies to $L(e,d)$.
\end{proof}

\begin{remark}
The capacity bound $K$ is not a mere technicality. Recall Brouwer's
invariance-of-domain theorem: a continuous injective map between open
subsets of $\mathbb R^n$ is necessarily an open map (its image is open); in
particular, no continuous injection of $\mathbb R^2$ into $\mathbb R^3$ can
have open image, so an exact space-filling reconstruction is topologically
impossible regardless of capacity. What Brouwer's theorem does \emph{not}
rule out is \emph{approximate} space-filling: an unconstrained sequence of
increasingly ill-conditioned (large-Lipschitz-constant) maps can approach
such behavior arbitrarily closely without ever achieving it, which would
erode the uniform lower bound $c$ as $K \to \infty$. In practice this is not
a loophole: the $\ell_2$ weight regularization and early stopping used when
training the autoencoder keep the realized Lipschitz constant controlled,
which is exactly why the dichotomy above is what we observe empirically.
\end{remark}

Together, Propositions 1 and 2 say that, at any fixed, reasonably regularized
network capacity, whether reconstruction error tends to $0$ or stays bounded
away from it is not a matter of degree: it is decided entirely by whether
the sampled data is confined to a $2$-dimensional subset of $\mathbb R^3$,
i.e.\ by whether $f_1, f_2, f_3$ satisfy an exact functional relation on the
region sampled. This is the theoretical content behind the empirical
detection rule used throughout this paper.

\section{Calibrating ``near zero'' with a permutation null}

An absolute error threshold is not meaningful across different triples of
scalars: they differ in units, scale, and marginal shape, and — since all
derived scalars are ultimately functions of the same three underlying
triangle parameters — some triples are mildly correlated even without an
exact relation. What is needed is a triple-specific baseline for ``no
relation.''

We construct a null sample $X'$ by independently permuting each of the three
columns of $X$. This exactly preserves each $f_i$'s own one-dimensional
marginal distribution while completely destroying the joint dependency
between the three: each row of $X'$ combines values from unrelated triangle
instances. Any real functional relation is annihilated by this shuffle.
Training the same autoencoder architecture on $X'$ (repeated a few times)
estimates the reconstruction error one should expect for three quantities
with these particular marginals but \emph{no} shared structure. Writing
$L_{\mathrm{real}}$ for the reconstruction risk on the real data and
$\mu_{\mathrm{null}}, \sigma_{\mathrm{null}}$ for the mean and standard
deviation of the reconstruction risk over the shuffled repeats, we compare
the two via the standard $z$-score
\[
  z \;=\; \frac{\mu_{\mathrm{null}} - L_{\mathrm{real}}}{\sigma_{\mathrm{null}}} ,
\]
i.e.\ the number of null standard deviations by which the real error falls
below the null mean, or equivalently via the ratio
$L_{\mathrm{real}}/\mu_{\mathrm{null}}$. Both
give a self-calibrated significance measure that is robust to the units and
distribution of whichever three scalars are being tested: a large $z$ (or a
small ratio) — real error far below the null — is the fingerprint of an
unexpected relation among exactly three quantities.

\section{The full pipeline}

\begin{enumerate}
  \item Enumerate candidate triples of derived scalar quantities.
  \item For each triple, sample random triangles once, train a 2-bottleneck
    autoencoder, and compare its held-out reconstruction error to a
    column-shuffled null (several repeats).
  \item Rank triples by how much better the real data compresses than its
    null counterpart (small ratio $=$ strong candidate).
\end{enumerate}

The autoencoder is used purely as a screening tool here: it tells us
\emph{whether} a triple of quantities is suspiciously dependent, not what
the relation is. The \texttt{verify\_euler\_relation} module (in
\texttt{triangle\_relations/discovery}) runs this exact pipeline on
$(R, r, d)$ as a worked example, confirming it correctly flags this triple —
Euler's relation — with no prior knowledge of the formula.

\end{document}
